**Proof.** Fix \(i\). For each \(u\in C_i\), let
\[
M(u):=\left\{t\in T:\; u\cdot t=\max_{r\in T} u\cdot r\right\}.
\]
Since \(T\) is finite, \(M(u)\neq\varnothing\). Also, by the definition of \(C_i\), the vertex \(s_i\) is the unique maximizer of \(u\cdot(\cdot)\) on \(S\).

We first prove the key claim:

> If \(u\in C_i\) and \(t\in M(u)\), then \(s_i+t\) is a lonely point.

Indeed, suppose
\[
s_i+t=s_j+t'
\]
for some \(s_j\in S\) and \(t'\in T\). Since \(t\in M(u)\), we have
\[
u\cdot t' \le u\cdot t.
\]
If \(j\neq i\), then \(u\in C_i\) implies
\[
u\cdot s_j<u\cdot s_i.
\]
Therefore
\[
u\cdot(s_j+t')
\le u\cdot s_j+u\cdot t
< u\cdot s_i+u\cdot t
= u\cdot(s_i+t),
\]
which contradicts \(s_i+t=s_j+t'\). Hence \(j=i\). Then
\[
s_i+t'=s_i+t,
\]
so \(t'=t\). Thus \(s_i+t\) has exactly one representation as \(s+t\) with \(s\in S\) and \(t\in T\), i.e. it is lonely.

Now let
\[
\ell_{s_i}=s_i+t_{s_i}
\]
be the unique lonely point associated to \(s_i\) given by Subproblem 1. By the claim, every \(t\in M(u)\) produces a lonely point \(s_i+t\). Since \(s_i\) has only one lonely point, we must have
\[
s_i+t=\ell_{s_i}=s_i+t_{s_i},
\]
and hence
\[
t=t_{s_i}.
\]

So every element of \(M(u)\) is equal to \(t_{s_i}\). Because \(M(u)\neq\varnothing\), it follows that
\[
M(u)=\{t_{s_i}\}
\qquad\text{for every }u\in C_i.
\]

Therefore the maximizing translate is independent of \(u\in C_i\). If we define
\[
t_i:=t_{s_i},
\]
then for every \(u\in C_i\),
\[
\operatorname*{arg\,max}_{t\in T} u\cdot t=\{t_i\}.
\]
Equivalently, \(t_i\) is the unique maximizer of \(u\cdot(\cdot)\) on \(T\) for every \(u\in C_i\). Moreover,
\[
\ell_{s_i}=s_i+t_i.
\]

This is exactly the desired statement. \(\square\)